\documentclass[10pt,a4paper]{amsart}
\usepackage[margin=19mm]{geometry}
\usepackage{amsmath,amssymb,mathtools}
\usepackage[scr=rsfs,cal=euler]{mathalfa}
\usepackage{xfrac}
\usepackage[unicode,hidelinks,bookmarks=false]{hyperref}
\newcommand{\ie}{i.e.\ }
\newcommand{\cf}{cf.\ }
\newcommand{\resp}{respectively\ }
\newcommand{\Subsection}{\subsection}
\providecommand{\nobreakdash}{\nobreak\hbox{-}}
\setlength{\emergencystretch}{2em}
\multlinegap0pt
\usepackage{bm}
\usepackage{bbm}
\usepackage{dsfont}
\usepackage{xspace}
\usepackage{cancel}
\usepackage{mathtools}
\usepackage{amsthm}
\makeatletter
\@ifundefined{equation}{\newtheorem{equation}{Equation}}{}
\@ifundefined{lemma}{\newtheorem{lemma}[equation]{Lemma}}{}
\@ifundefined{remark}{\newtheorem{remark}[equation]{Remark}}{}
\makeatother
\title[Homogeneous spaces with two equivalent isotropy summands dev]{Homogeneous spaces with two equivalent isotropy summands develop positive Ricci curvature under Ricci flow}
\author{Eric Cochran, Arseny Mingajev, Lawrence Mouill\'e,, Nazia Valiyakath}
\date{Source: arXiv:2608.02840v1 (CC BY 4.0)}
\begin{document}
\maketitle

\noindent\emph{Source and licence.} Homogeneous spaces with two equivalent isotropy summands develop positive Ricci curvature under Ricci flow. Version arXiv:2608.02840v1 (CC BY 4.0), distributed under the Creative Commons Attribution 4.0 International licence, as recorded in the arXiv licence metadata for this exact version and shown on the abstract page https://arxiv.org/abs/2608.02840v1. Full rights notice: licenses/arxiv-2608.02840-CC-BY-4.0.txt.\par\smallskip

\noindent\textit{Editorial scope.} Counted unit: the Remark whose source label is \texttt{rem:compact} in the article (source0.tex lines 892--897, source label \texttt{rem:compact}), delivered together with the article's own complete original proof of it (source0.tex lines 974--1056, one \texttt{proof} environment of 83 lines). No mathematics is written, repaired or re-derived: statement and proof are the article's own text line for line. The proof is detached from the statement in the source: the article states the result at lines 892--897 and prints its proof at lines 974--1056, and the proof environment's own header names the statement, so the association is the article's own. Required context retained uncounted: eq:sys\_xy (lines 703--712); curve1xy (lines 719--724); curve2xy (lines 719--724); curve3xy (lines 719--724); curve4xy (lines 719--724); lem:Pinc (lines 800--805). Every definition, display and label of those lines is retained unchanged. Omitted from this package and not redistributed: 1--702, 713--718, 725--799, 806--891, 898--973, 1057--1455. Nothing inside a retained excerpt is omitted or altered: every retained body is delivered exactly as the article's lines are written, so no omission receipt of any kind accompanies this package. Citations: the retained excerpts carry no citation command at all, so no bibliography entry of the article is transcribed and no reference is redistributed. Reference adaptation: none was needed and none was made; every reference inside the delivered unit resolves inside it, so no unresolved reference or citation is produced. Printed slips kept literal: no printed word, symbol, hypothesis or number of the counted statement or of its proof was altered. Layout and class: the article's own class is \texttt{amsart} with packages including inputenc, amsmath, amsfonts, amssymb, amsthm, fontenc, etoolbox, parskip, hyperref, tikz, todonotes, enumerate, float, amsaddr; the wrapper is the lane's pinned \texttt{amsart} editorial wrapper at 10pt on A4, which restates the theorem-like environments the retained text needs and adds the article's own macro definitions that the retained text needs (none), with extra wrapper packages (bm, bbm, dsfont, xspace, cancel, mathtools). The delivered numbering is the unit's own. Assets: no retained span contains a figure environment, a graphics-inclusion command, a PDF-page inclusion, a table or a TikZ picture; no image file is copied into this package, the package declares no asset dependency and no image file is redistributed with it. Licence: version arXiv:2608.02840v1 of this article is distributed under the Creative Commons Attribution 4.0 International licence, as recorded in the arXiv licence metadata for this exact version and shown on the abstract page https://arxiv.org/abs/2608.02840v1. A full rights notice is delivered at licenses/arxiv-2608.02840-CC-BY-4.0.txt. Characters: every retained excerpt is pure ASCII, so the delivered engine, XeLaTeX with its default Latin Modern text font, reports no missing character.
    \begin{equation}\label{eq:sys_xy}
    \begin{aligned} 
        &
        x'(t) = -6x\left(x^3 + 3x^2 y + 4x + 6y\right), \\
        &
        y'(t) = -6x^3 y - 18x^2 y^2 + 6x^2 + 18x y + 36y^2 - 6. \\
        &
        2xy+6y^2\geq 1, \quad y>0  .
    \end{aligned}
    \end{equation}

    \begin{enumerate}
        \item \label{curve1xy} $2xy + 6y^2 = 1$ (with $y>0$),
        \item \label{curve2xy} $x=0$,
        \item \label{curve3xy} $3x^2 + 8xy + 2 = 0$ (with $x<0$), and %$3x^2 + 8xy + 2 = 0$,
        \item \label{curve4xy} $2x^2 + 6xy + 3 = 0$ (with $x<0$). %$2x^2 + 6xy + 3 = 0$.
    \end{enumerate}

    \begin{lemma}
        \label{lem:Pinc}
        For any $(x,y)$ lying in the region where $2xy+6y^2>1$ and $y>0$, we have $P'(t)>0$ for all $t$ except on the two invariant curves $x=0$ and $2x^2+6xy+3=0$ (with $x<0$), on which $P'(t)$ is zero. 
        Additionally, these two curves are contained in the region where $P>0$.
        Consequently, the region where $P>0$ is forward-invariant.
    \end{lemma}

    \begin{remark}\label{rem:compact}
        Note that by Lemma \ref{lem:Pinc}, System \eqref{eq:sys_xy} has no cyclic solutions, as such a solution would necessarily have $P'(t) \equiv 0$, and this equality only holds on the unbounded curves \eqref{curve2xy} and \eqref{curve4xy}.
        Consequently by the Poincar\'e-Bendixson theorem, because the fixed points of the system consist of a source at $(x,y)=(-\sqrt{2},\frac{\sqrt{2}}{2})$ and saddle points at $(x,y)=(-\sqrt{6},\frac{5\sqrt{6}}{12})$ and $(0,\frac{\sqrt{6}}{6})$, the only solutions which are contained in a compact region are the fixed points and solutions on the curves which serve as the stable manifolds for the saddle points.
        As stated above, these stable manifolds lie on curves \eqref{curve1xy} or \eqref{curve3xy}.
        Therefore, because we only need to work with initial conditions which do not lie on these invariant curves, we may assume all solutions analyzed below are unbounded.
    \end{remark}

    \begin{proof}
        If the initial condition lies in region $F$ (where $x>0$), then $x(t)>0$ for all $t$, and thus $x'(t) = -6x\left(x^3 + 3x^2 y + 4x + 6y\right)<0$ for all $t$.
        If instead the initial condition lies in region $C$, then it can be verified with Maple's IsEmpty function that $x'(t)>0$ for all $t$.
        Therefore, $|x(t)|$ is monotonically decreasing for all $t$ in both cases.
        
        Next, we show that $y(t)$ eventually increases monotonically.
        By Remark \ref{rem:compact}, we may assume that our solution $(x(t),y(t))$ is unbounded.
        Since $|x|$ is bounded, we have that $|y|$ is unbounded.  
        Because $y>0$ by assumption, the solution must satisfy $y'(t) > 0$ for some $t$.  
        Using Maple's IsEmpty function, we can verify if $y>2$, then $y''(t)>0$.
        Thus the region where $y>2$ and $y'(t)>0$ is forward-invariant, and because $y$ is unbounded, our solution must eventually enter and remain in this region.
        Hence, $y(t)$ increases monotonically for all sufficiently large values of $t$, and because $y$ is unbounded, we have $y\to \infty$ as $t$ increases.
        
        

        % Now we show that $x\to 0$ as $t\to\infty$.
        % The region where $y'(t)>0$ has two connected components which are separated by the line $18y+10x=0$; see Figure \ref{fig:y_prime_pos_sep_by_line}.
        % We observe that, because $|x|$ is bounded and $y$ is positive and unbounded, our solution cannot be contained in the region where $18y+10x<0$ and hence must eventually enter and remain in the region where $y>2$, $y'(t)>0$, and $18y+10x>0$.
        % Furthermore, it can be seen directly in \eqref{eq:sys_xy} that $x'(t) < 0$ if $x > 0$, and it can be verified with Maple that $x'(t) > 0$ if $x < 0$, $y>2$, $y'(t)>0$, and $18y+10x>0$.  
        % Since $x=0$ is an invariant curve, the sign of $x$ for our solution will not change, and therefore $|x(t)|$ eventually decreases monotonically.

        % \begin{figure}
        %     \centering
        %     \includegraphics[width=0.6\linewidth]{divided_region.png}
        %     \caption{The region $y'(t)>0$ separated by the line $18y+10x=0$.}
        %     \label{fig:y_prime_pos_sep_by_line}
        % \end{figure}

        Consequently, we may view $y$ as an invertible function of $x$ for sufficiently large values of $t$.
        Thus, considering $\frac{dy}{dx} = \frac{y'(t)}{x'(t)}$, we have by \eqref{eq:sys_xy} that 
        \[ 
            \frac{dy}{dx} = \frac{-6x^3 y - 18x^2 y^2 + 6x^2 + 18x y + 36y^2 - 6}{-6x\left(x^3 + 3x^2 y + 4x + 6y\right)}. 
        \]
        % Since $|x(t)|$ is bounded and eventually monotonically decreasing while $y(t)$ is unbounded and eventually monotonically increasing, we can consider the limit of $\frac{dy}{dx} \cdot \frac{1}{y}$ as $y\to \infty$, where we treat $x$ as a function of $y$.
        Using that $|x|$ is bounded, we get that
        \begin{equation}\label{eqn A}
             \lim_{y \to \infty}\frac{dy}{dx} \cdot \frac{1}{y} = \lim_{y \to \infty}\frac{x^2-2}{x\left(x^2+2\right)}. 
        \end{equation}
        Because $|x(t)|$ is bounded and monotonic, it must converge to some finite value as $t$ increases (equiv. $y\to \infty$).
        For the sake of contradiction, assume this limit is non-zero.
        % Thus as $y\to \infty$, 
        % This is equivalent to letting $y \to \infty$, and we see in this case that the above limit does exist in this case (let's call the limit $k$) since that would imply the expression above is (eventually) monotonic and bounded. 
        Then the limit \eqref{eqn A} converges to some finite value, and thus there exists a positive number $M$ such that, for sufficiently large values of $t$, we have
        \[
            \left | \frac{dy}{dx} \cdot \frac{1}{y} \right | < M.
        \]
        % Hence, the growth rate of $|\frac{dy}{dx}|$ is eventually dominated by a linear function.  
        Because $|x(t)|$ is bounded, the above inequality implies that $|y(t)|$ must also be bounded, which is a contradiction.
        Thus, we have $x\to 0$ as $t$ increases.
        %
        % \todo[inline]{Line separating new (above) from old (below).}
        %
        % Hence, we can assume $C$ is contained in the connected component of the region $y'>0$ above the line $18y+10x=0$. 
        % A computation in \textsc{MAPLE} shows that the region for which $y'>0$, $18y+10x>0$, and $y''<0$ is empty, meaning the connected component lying above this line is forward-invariant.
        % Furthermore, $x' > 0$ if $x < 0$ and $x' < 0$ if $x > 0$ in this component.  
        % Since $x=0$ is an invariant curve, note that our integral curve $C$ will always have either $x<0$ and $x'>0$ or $x>0$ and $x'<0$.  In other words, $|x(t)|$ eventually decreases monotonically and $y(t)$ eventually increases monotonically.  
        %
        %
        % Eventual monotonicity of $x(t)$ allows us to view $y$ as a function of $x$ for large values of $t$.  
        % We can thus consider the derivative $\frac{dy}{dx} = \frac{dy/dt}{dx/dt}$.
        % Observe that this quantity is always defined since $\frac{dx}{dt} \neq 0$.  
        % Furthermore, by \eqref{eq:sys_xy}, we have
        % \[ \frac{dy}{dx} = \frac{-6x^3 y - 18x^2 y^2 + 6x^2 + 18x y + 36y^2 - 6}{-6x\left(x^3 + 3x^2 y + 4x + 6y\right)}. \]
        %
        % Since $|x(t)|$ is bounded and eventually monotonically decreasing while $y(t)$ is unbounded and eventually monotonically increasing, we can consider the limit of $\frac{dy}{dx} \cdot \frac{1}{y}$ as $y\to \infty$, where we treat $x$ as a function of $y$.
        % Using that $|x|$ is bounded, we have
        % \begin{equation}\label{eqn A}
        %      \lim_{y \to \infty}\frac{dy}{dx} \cdot \frac{1}{y} = \lim_{y \to \infty}\frac{-18x^2+36}{-6x\left(3x^2+6\right)}. 
        % \end{equation}
        %
        %
        % Because $x(t)$ is bounded and eventually monotonic, it must converge to some finite value as $t \to \infty$ (equiv. $y\to \infty$).
        % For the sake of contradiction, assume this limit is non-zero.
        % % Thus as $y\to \infty$, 
        % % This is equivalent to letting $y \to \infty$, and we see in this case that the above limit does exist in this case (let's call the limit $k$) since that would imply the expression above is (eventually) monotonic and bounded. 
        % Then the limit \eqref{eqn A} converges to some $k$, and thus given $\epsilon > 0$, for sufficiently large values of $t$, we have that
        % \[
        %     \left | \frac{dy}{dx} \right | < |k|y+\epsilon
        % \]
        % Hence, the growth rate of $|\frac{dy}{dx}|$ is eventually dominated by a linear function.  
        % This means that as $y$ tends to infinity, $x$ must also tend to infinity. This means that if $|x(t)|$ is bounded, $|y(t)|$ must also be bounded, contradicting Lemma \ref{lem:compact}.
        % Therefore, $|x(t)|$ must go to zero monotonically as $t \to \infty$.
    \end{proof}

\end{document}
